% Module 1: Bits to Qubits
% Estimated slides: 35 | Duration: ~2 hours
% Key topics: Classical bits, qubit definition, superposition, Bloch sphere, measurement,
%             multi-qubit state space, quantum advantage intuition
% Labs: Qiskit -- creating and measuring qubits, Bloch sphere visualization

%%%%%%%%%%%%%%%%%%%%%%%%%%%%%%%%%%%%%%%%%%%%%%%%%%%%%%%%%%
\begin{frame}[fragile]\frametitle{}
\begin{center}
{\Large Module 1: From Bits to Qubits}
\end{center}
\end{frame}

%%%%%%%%%%%%%%%%%%%%%%%%%%%%%%%%%%%%%%%%%%%%%%%%%%%%%%%%%%%
\begin{frame}[fragile]\frametitle{Classical Bits: The Building Block of Computing}
\begin{itemize}
\item Every classical computer stores information as bits: 0 or 1
\item Physically realized as: voltage levels (low/high), magnetic domains, optical pulses
\item All classical computation: manipulating sequences of 0s and 1s
\item A byte = 8 bits, represents 256 possible values
\item An $n$-bit register holds exactly ONE value at a time out of $2^n$ possibilities
\end{itemize}
% Suggested diagram: Physical realizations of bits (voltage, hard drive, CD)
\end{frame}

%%%%%%%%%%%%%%%%%%%%%%%%%%%%%%%%%%%%%%%%%%%%%%%%%%%%%%%%%%%
\begin{frame}[fragile]\frametitle{Bits to Qubits}
\begin{itemize}
\item Classical bit: state is 0 or 1
\item Qubit: state is a vector in 2-dimensional complex vector space
\item Superposition: continuous states, like flipping coin
\item When observed, qubit collapses to 0 or 1
\item Before observation: probabilistically in middle states
\item For $n$ qubits: $2^n$ possible collapsing values
\item Calculation happens on whole continuous function simultaneously
\end{itemize}
\end{frame}

%%%%%%%%%%%%%%%%%%%%%%%%%%%%%%%%%%%%%%%%%%%%%%%%%%%%%%%%%%%
\begin{frame}[fragile]\frametitle{The Qubit: A Quantum Bit}
\begin{itemize}
\item A qubit is the quantum analogue of a classical bit
\item Like a classical bit, when measured it gives either 0 or 1
\item Unlike a classical bit, before measurement it can be in a superposition of both
\item Physically realized as: electron spin, photon polarization, superconducting circuit state
\item The measurement outcome is probabilistic, not deterministic
\end{itemize}
% Suggested diagram: Coin analogy -- spinning coin (superposition) vs flat coin (classical bit)
\end{frame}

%%%%%%%%%%%%%%%%%%%%%%%%%%%%%%%%%%%%%%%%%%%%%%%%%%%%%%%%%%%
\begin{frame}[fragile]\frametitle{Quantum State of a Qubit}
\begin{itemize}
\item State is a vector of unit length in 2D complex vector space
\item Called state space
\item For example: $\begin{bmatrix} 1 \\ 0 \end{bmatrix}$ or $\begin{bmatrix} 0.6 \\ 0.8 \end{bmatrix}$
\item General form: $\alpha|0\rangle + \beta|1\rangle$ where $\alpha, \beta$ are complex numbers
\item Ket notation: $|0\rangle$ and $|1\rangle$ represent basis states
\item Contains all information about quantum system
\end{itemize}
\end{frame}

%%%%%%%%%%%%%%%%%%%%%%%%%%%%%%%%%%%%%%%%%%%%%%%%%%%%%%%%%%%
\begin{frame}[fragile]\frametitle{Superposition: The Coin Analogy}
\begin{itemize}
\item Imagine a coin:
  \begin{itemize}
  \item Classical: coin flat on table --- heads (0) or tails (1), definite state
  \item Quantum: coin spinning in the air --- neither heads nor tails until it lands
  \end{itemize}
\item A qubit in superposition has a \textit{probability amplitude} for 0 and for 1
\item When we measure, it ``lands'' on 0 or 1 with probabilities set by the amplitudes
\item The probabilities must sum to 1 (certainty)
\item Key: the probabilities are determined before measurement but the outcome is random
\end{itemize}
% Suggested diagram: Spinning coin -> landing = measurement; probability pie chart
\end{frame}

%%%%%%%%%%%%%%%%%%%%%%%%%%%%%%%%%%%%%%%%%%%%%%%%%%%%%%%%%%%
\begin{frame}[fragile]\frametitle{Superposition Explained}
\begin{itemize}
\item Qubit exists in multiple states at once until measured
\item Calculation processes all superposed values together
\item Output calculated for all continuous values simultaneously
\item Gate operations manipulate probability amplitudes
\item Clever gates increase probability of desired outcome
\item When measured, most likely to collapse to solution
\item Enables massive parallelism in computation
\end{itemize}
\end{frame}

%%%%%%%%%%%%%%%%%%%%%%%%%%%%%%%%%%%%%%%%%%%%%%%%%%%%%%%%%%%
\begin{frame}[fragile]\frametitle{Probability Amplitudes (No Heavy Math)}
\begin{itemize}
\item A qubit state is described by two numbers $\alpha$ and $\beta$ (complex amplitudes)
\item The probability of measuring 0 is $|\alpha|^2$; probability of measuring 1 is $|\beta|^2$
\item They must satisfy: $|\alpha|^2 + |\beta|^2 = 1$ (probabilities sum to 1)
\item Example: $\alpha = \beta = 1/\sqrt{2}$ gives 50\% chance of 0, 50\% chance of 1
\item This is the \textit{uniform superposition} state (like a fair coin)
\end{itemize}
% Suggested diagram: Circle with two sectors labelled |alpha|^2 and |beta|^2
\end{frame}

%%%%%%%%%%%%%%%%%%%%%%%%%%%%%%%%%%%%%%%%%%%%%%%%%%%%%%%%%%%
\begin{frame}[fragile]\frametitle{The Bloch Sphere: Visualizing a Qubit}
\begin{itemize}
\item Every possible state of a qubit maps to a point on the surface of a sphere
\item North pole = state $|0\rangle$ (classical 0)
\item South pole = state $|1\rangle$ (classical 1)
\item Equator = equal superposition (50/50 chance of 0 or 1)
\item Any point in between = some superposition
\item The angle from the north pole determines the probability ratio
\item The ``longitude'' angle encodes the relative phase (important for interference)
\end{itemize}
% Suggested diagram: Bloch sphere with labeled poles, equator, and an example state vector
\end{frame}

%%%%%%%%%%%%%%%%%%%%%%%%%%%%%%%%%%%%%%%%%%%%%%%%%%%%%%%%%%%
\begin{frame}[fragile]\frametitle{Quantum Measurement: Collapse}
\begin{itemize}
\item Before measurement: qubit is in superposition (can be anywhere on Bloch sphere)
\item Measurement: we ask ``is this qubit 0 or 1?'' and get a definite answer
\item After measurement: the qubit collapses to either $|0\rangle$ or $|1\rangle$
\item The superposition is destroyed; the qubit is now a classical bit
\item Every subsequent measurement gives the same result (no superposition left)
\item This is why we run quantum circuits many times to get probability distributions
\end{itemize}
% Suggested diagram: Before-measurement Bloch sphere state -> arrow -> measurement outcome -> collapsed state
\end{frame}

%%%%%%%%%%%%%%%%%%%%%%%%%%%%%%%%%%%%%%%%%%%%%%%%%%%%%%%%%%%
\begin{frame}[fragile]\frametitle{Measuring a Qubit}
\begin{itemize}
\item Measurement in computational basis extracts classical information
\item State $\alpha|0\rangle + \beta|1\rangle$ gives outcome 0 with probability $|\alpha|^2$
\item Gives outcome 1 with probability $|\beta|^2$
\item Measurement disturbs the state (no longer superposition)
\item Posterior state: $|0\rangle$ if outcome 0, $|1\rangle$ if outcome 1
\item Amplitudes $\alpha, \beta$ are hidden information after measurement
\item Cannot extract infinite classical information from qubit
\end{itemize}
\end{frame}

%%%%%%%%%%%%%%%%%%%%%%%%%%%%%%%%%%%%%%%%%%%%%%%%%%%%%%%%%%%
\begin{frame}[fragile]\frametitle{Measurement Probabilities}
\begin{itemize}
\item Probabilities must sum to 1: $|\alpha|^2 + |\beta|^2 = 1$
\item This is exactly the normalization condition
\item Origin: measurement probabilities are physical requirements
\item Measurement is how we extract computation results
\item Circuit notation: semicircle symbol with double wire output
\item Classical bit $m$ denotes measurement result (0 or 1)
\item After measurement, quantum information is classical
\end{itemize}
\end{frame}

%%%%%%%%%%%%%%%%%%%%%%%%%%%%%%%%%%%%%%%%%%%%%%%%%%%%%%%%%%%
\begin{frame}[fragile]\frametitle{Why Can't We Just Look at a Qubit?}
\begin{itemize}
\item Observation disturbs quantum systems --- this is fundamental, not a technical limitation
\item You cannot check a qubit's state without collapsing it
\item This is related to the Heisenberg Uncertainty Principle
\item This has practical consequences for quantum algorithms:
  \begin{itemize}
  \item We cannot read out partial results mid-computation
  \item We must design algorithms so the right answer has high probability at the end
  \end{itemize}
\item Quantum debugging is harder than classical debugging!
\end{itemize}
\end{frame}

%%%%%%%%%%%%%%%%%%%%%%%%%%%%%%%%%%%%%%%%%%%%%%%%%%%%%%%%%%%
\begin{frame}[fragile]\frametitle{Multiple Qubits: Exponential State Space}
\begin{itemize}
\item 1 qubit: 2 possible states (0 or 1)
\item 2 qubits: 4 possible states (00, 01, 10, 11)
\item 3 qubits: 8 possible states
\item $n$ qubits: $2^n$ possible states
\item A quantum computer with $n$ qubits can be in a superposition of ALL $2^n$ states
\item 50 qubits: $2^{50} \approx 10^{15}$ states simultaneously
\item 300 qubits: more states than atoms in the observable universe
\end{itemize}
% Suggested diagram: Table showing n qubits, 2^n states, and practical meaning
\end{frame}

%%%%%%%%%%%%%%%%%%%%%%%%%%%%%%%%%%%%%%%%%%%%%%%%%%%%%%%%%%%
\begin{frame}[fragile]\frametitle{Quantum Parallelism: The Key Advantage}
\begin{itemize}
\item A quantum computer can apply an operation to ALL $2^n$ states simultaneously
\item This is not the same as running $2^n$ parallel classical computers
\item The trick: we can only extract $n$ bits of information on measurement
\item Quantum algorithms must cleverly funnel amplitudes towards the right answer
\item The art of quantum algorithms = interference engineering
\item Constructive interference: amplify the correct answer
\item Destructive interference: suppress the wrong answers
\end{itemize}
% Suggested diagram: Wave interference analogy -- constructive (peaks add) and destructive (cancel)
\end{frame}

%%%%%%%%%%%%%%%%%%%%%%%%%%%%%%%%%%%%%%%%%%%%%%%%%%%%%%%%%%%
\begin{frame}[fragile]\frametitle{Comparing Classical and Quantum Bits}
\begin{itemize}
\item \textbf{Classical bit}: 0 or 1, deterministic, cheap to copy, read without disturbance
\item \textbf{Qubit}: probabilistic mixture before measurement, collapses on reading, cannot be copied (no-cloning)
\item Both collapse to 0 or 1 when read --- the difference is what happens in between
\item Quantum states live in a much larger mathematical space (complex amplitudes)
\item This extra ``richness'' is the source of quantum computational power
\end{itemize}
% Suggested diagram: Comparison table Classical vs Quantum (bit, state, copy, read, parallelism)
\end{frame}

%%%%%%%%%%%%%%%%%%%%%%%%%%%%%%%%%%%%%%%%%%%%%%%%%%%%%%%%%%%
\begin{frame}[fragile]\frametitle{Physical Realizations: What IS a Qubit?}
\begin{itemize}
\item \textbf{Electron spin}: spin up = $|0\rangle$, spin down = $|1\rangle$
\item \textbf{Photon polarization}: horizontal = $|0\rangle$, vertical = $|1\rangle$
\item \textbf{Superconducting circuit}: ground energy state = $|0\rangle$, excited = $|1\rangle$
\item \textbf{Trapped ion}: two energy levels of an ion atom
\item All of these behave as two-level quantum systems
\item The abstraction (qubit) is the same regardless of physical implementation
\end{itemize}
% Suggested diagram: Icons for each physical system
\end{frame}

%%%%%%%%%%%%%%%%%%%%%%%%%%%%%%%%%%%%%%%%%%%%%%%%%%%%%%%%%%%
\begin{frame}[fragile]\frametitle{Lab Preview: Your First Qubit in Qiskit}
\begin{itemize}
\item We will create a single qubit in the $|0\rangle$ state
\item Apply a Hadamard gate to put it in equal superposition
\item Measure it 1000 times and plot the histogram of results
\item Expected result: approximately 500 zeros and 500 ones
\item This demonstrates superposition and probabilistic measurement
\end{itemize}
\begin{verbatim}
from qiskit import QuantumCircuit
from qiskit_aer import AerSimulator
qc = QuantumCircuit(1, 1)
qc.h(0)           # Hadamard: create superposition
qc.measure(0, 0)  # Measure qubit 0
# Run on simulator and plot histogram
\end{verbatim}
% Suggested diagram: Circuit diagram showing |0> -> H -> Measure
\end{frame}

%%%%%%%%%%%%%%%%%%%%%%%%%%%%%%%%%%%%%%%%%%%%%%%%%%%%%%%%%%%
\begin{frame}[fragile]\frametitle{Lab: Bloch Sphere Visualization}
\begin{itemize}
\item Qiskit can visualize qubit states on the Bloch sphere
\item We will plot several common states: $|0\rangle$, $|1\rangle$, $|+\rangle$, $|-\rangle$
\item This builds geometric intuition for quantum states
\end{itemize}
\begin{verbatim}
from qiskit.visualization import plot_bloch_multivector
from qiskit.quantum_info import Statevector
sv = Statevector.from_label('+')
plot_bloch_multivector(sv)
\end{verbatim}
% Suggested diagram: Bloch sphere screenshot from Qiskit
\end{frame}

%%%%%%%%%%%%%%%%%%%%%%%%%%%%%%%%%%%%%%%%%%%%%%%%%%%%%%%%%%%
\begin{frame}[fragile]\frametitle{Qubit States: Key Reference Points}
\begin{itemize}
\item $|0\rangle$ -- north pole, always measures 0
\item $|1\rangle$ -- south pole, always measures 1
\item $|+\rangle = (|0\rangle + |1\rangle)/\sqrt{2}$ -- equator, 50/50 in Z basis
\item $|-\rangle = (|0\rangle - |1\rangle)/\sqrt{2}$ -- equator, 50/50 in Z basis but opposite phase
\item $|i\rangle = (|0\rangle + i|1\rangle)/\sqrt{2}$ -- equator, Y axis
\item The phase difference between $|+\rangle$ and $|-\rangle$ matters for interference
\end{itemize}
% Suggested diagram: Bloch sphere with all 6 cardinal states labeled
\end{frame}

%%%%%%%%%%%%%%%%%%%%%%%%%%%%%%%%%%%%%%%%%%%%%%%%%%%%%%%%%%%
\begin{frame}[fragile]\frametitle{Computational Basis States}
\begin{itemize}
\item Two special states correspond to classical 0 and 1
\item $|0\rangle := \begin{bmatrix} 1 \\ 0 \end{bmatrix}$ (computational basis state for 0)
\item $|1\rangle := \begin{bmatrix} 0 \\ 1 \end{bmatrix}$ (computational basis state for 1)
\item These are orthonormal vectors
\item Serve as foundation for general quantum states
\item Any qubit state is linear combination of these
\end{itemize}
\end{frame}

%%%%%%%%%%%%%%%%%%%%%%%%%%%%%%%%%%%%%%%%%%%%%%%%%%%%%%%%%%%
\begin{frame}[fragile]\frametitle{General States of a Qubit}
\begin{itemize}
\item Beyond $|0\rangle$ and $|1\rangle$, many more states possible
\item Example: $0.6|0\rangle + 0.8|1\rangle$
\item Superposition: linear combination of basis states
\item Amplitude: coefficient for particular state
\item In general: $\alpha|0\rangle + \beta|1\rangle$ where $\alpha, \beta \in \mathbb{C}$
\item These extra states give quantum computers power beyond classical
\end{itemize}
\end{frame}

%%%%%%%%%%%%%%%%%%%%%%%%%%%%%%%%%%%%%%%%%%%%%%%%%%%%%%%%%%%
\begin{frame}[fragile]\frametitle{Normalization Constraint}
\begin{itemize}
\item Sum of squares of amplitudes must equal 1
\item For state $\alpha|0\rangle + \beta|1\rangle$: $|\alpha|^2 + |\beta|^2 = 1$
\item Example: $(0.6)^2 + (0.8)^2 = 0.36 + 0.64 = 1$ \checkmark
\item Ensures quantum state is unit vector
\item Related to measurement probabilities summing to 1
\item Fundamental requirement for valid quantum states
\end{itemize}
\end{frame}

%%%%%%%%%%%%%%%%%%%%%%%%%%%%%%%%%%%%%%%%%%%%%%%%%%%%%%%%%%%
\begin{frame}[fragile]\frametitle{Quantum States: The Water Glass Analogy}
\begin{itemize}
\item Imagine filling a glass with colored water as a metaphor for probability amplitudes
\item The fraction of water in the 0-half = probability of measuring 0
\item You can tilt the glass (rotate the state) without spilling (unitary operation)
\item When you set the glass down flat (measure), it snaps to one side
\item The snap is random but weighted by how much water was on each side
\item Gates = controlled tilting; measurement = letting it settle
\end{itemize}
\end{frame}

%%%%%%%%%%%%%%%%%%%%%%%%%%%%%%%%%%%%%%%%%%%%%%%%%%%%%%%%%%%
\begin{frame}[fragile]\frametitle{From One Qubit to Quantum Registers}
\begin{itemize}
\item A quantum register is $n$ qubits working together
\item The state of $n$ qubits lives in a $2^n$-dimensional space
\item Qiskit QuantumCircuit(n) creates an n-qubit register
\item Qubits in a register can be independent or entangled (correlated)
\item Entanglement is the key to quantum advantage over independent qubits
\item We will explore entanglement in Module 2
\end{itemize}
% Suggested diagram: Single qubit vs 2-qubit vs n-qubit register
\end{frame}

%%%%%%%%%%%%%%%%%%%%%%%%%%%%%%%%%%%%%%%%%%%%%%%%%%%%%%%%%%%
\begin{frame}[fragile]\frametitle{Why Classical Simulation of Qubits is Exponentially Hard}
\begin{itemize}
\item To simulate $n$ qubits classically, you need $2^n$ complex numbers
\item 50 qubits: $2^{50}$ = 1 quadrillion amplitudes = ~8 petabytes of memory
\item This is why quantum computers are hard to simulate and why they matter
\item Current classical supercomputers can simulate up to ~50-55 qubits
\item Beyond that, only a quantum computer can efficiently handle the state
\item This is the core of ``quantum computational advantage''
\end{itemize}
% Suggested diagram: Memory requirement curve (exponential) vs qubit count
\end{frame}

%%%%%%%%%%%%%%%%%%%%%%%%%%%%%%%%%%%%%%%%%%%%%%%%%%%%%%%%%%%
\begin{frame}[fragile]\frametitle{Exercise: Qubit Probability Calculations}
\begin{itemize}
\item A qubit has amplitudes $\alpha = \sqrt{3}/2$ for $|0\rangle$ and $\beta = 1/2$ for $|1\rangle$
\begin{enumerate}
  \item What is the probability of measuring 0? (Answer: 3/4)
  \item What is the probability of measuring 1? (Answer: 1/4)
  \item Do they sum to 1? Verify: $(3/4) + (1/4) = 1$
  \item After measuring 0, what is the qubit's state? (Answer: $|0\rangle$, collapsed)
\end{enumerate}
\item \textbf{Extended}: Write a Qiskit circuit that produces these probabilities using Ry gate
\end{itemize}
\end{frame}

%%%%%%%%%%%%%%%%%%%%%%%%%%%%%%%%%%%%%%%%%%%%%%%%%%%%%%%%%%%
\begin{frame}[fragile]\frametitle{Module 1 Summary}
\begin{itemize}
\item Classical bits are definite (0 or 1); qubits can be in superposition
\item Superposition = probability amplitudes for both 0 and 1 simultaneously
\item Measurement collapses the qubit to a definite value, destroying superposition
\item The Bloch sphere is the standard visual for a single qubit state
\item $n$ qubits span $2^n$ states simultaneously, enabling quantum parallelism
\item Quantum advantage comes from interference: amplifying right, suppressing wrong
\item Next: The fundamental rules that govern quantum systems
\end{itemize}
\end{frame}
